\selectlanguage{spanish}

\renewcommand{\articuloidiomaxml}{es}
\renewcommand{\articulotitulo}{Editorial}
\renewcommand{\articuloautores}{Juan Carlos García Buitrago}
\renewcommand{\articulotipo}{EDITORIAL}
\renewcommand{\articulorevista}{Bioeconomía y territorio}
\renewcommand{\articuloissne}{2500-7211}
\renewcommand{\articulovolumen}{9}
\renewcommand{\articulonumero}{1}
\renewcommand{\articulofecha}{ 2026}
\renewcommand{\articulopagina}{1}

\setcounter{page}{1}
\thispagestyle{firstpage}

% --- TÍTULO A ANCHO COMPLETO ---
\noindent\begin{minipage}[t]{\dimexpr\textwidth+\marginparsep+\marginparwidth\relax}%
{\sffamily\bfseries\Large\articulotitulo\par}%
\end{minipage}\par
\vspace{3mm}%
\renewcommand{\articulotitulotrans}{Editor's note}
\noindent\begin{minipage}[t]{\dimexpr\textwidth+\marginparsep+\marginparwidth\relax}%
{\sffamily\itshape\normalsize\articulotitulotrans\par}%
\end{minipage}\par
\vspace{4mm}%

% --- BLOQUE LATERAL DE METADATOS ---
\begin{textblock*}{68mm}(\gbbloquex,92mm)
{\setlength{\leftskip}{0pt}\setlength{\parindent}{0pt}%
\noindent\begin{minipage}[t]{68mm}
\sffamily\small\setlength{\parskip}{1pt}
{\bfseries\color{azulrevista}Juan Carlos García Buitrago}\par
\vspace{6pt}
{\bfseries\color{azulrevista!70}Derechos de autor\'ia:}\par
\textcopyright\ 2026 García Buitrago, J. Publicado por Bioeconomía y territorio. Se permite el uso, distribución y reproducción sin fines comerciales, siempre que se cite la obra original y las obras derivadas se distribuyan bajo la misma licencia. (\href{https://creativecommons.org/licenses/by-nc-sa/4.0/}{https://creativecommons.org/licenses/by-nc-sa/4.0/})\par\vspace{6pt}
{\bfseries\color{azulrevista!70}Publicación:} 2026, vol.~9, n\'um.~1, pp.~1--1\par\vspace{6pt}
{\bfseries\color{azulrevista!70}URL:} {\scriptsize\texttt{\href{https://www.dropbox.com/scl/fi/rom8ermbyet1p606ms6y0/l-01-EDITORIAL.pdf?rlkey=x341kromqoa6a1fyx499z5yj2&st=mknhry7g&dl=1}{https://www.dropbox.com/scl/fi/rom8ermbyet1p606ms6y0/l-01-EDITORIAL.pdf?rlkey=x341kromqoa6a1fyx499z5yj2&st=mknhry7g&dl=1}}}\par\vspace{6pt}
{\bfseries\color{azulrevista!70}Licencia:} CC BY-NC-SA 4.0\par
{\scriptsize\texttt{\href{https://creativecommons.org/licenses/by-nc-sa/4.0/}{https://creativecommons.org/licenses/by-nc-sa/4.0/}}}\par
\vspace{6pt}
\vspace{6pt}
{\bfseries\color{azulrevista!70}Compartir:}\par\vspace{3pt}
{\large
\href{https://bsky.app/intent/compose?text=https%3A%2F%2Fwww.dropbox.com%2Fscl%2Ffi%2From8ermbyet1p606ms6y0%2Fl-01-EDITORIAL.pdf%3Frlkey%3Dx341kromqoa6a1fyx499z5yj2%26st%3Dmknhry7g%26dl%3D1}{\includesvg[height=0.9em]{/home/alberto/.gbpublisher/fonts/bluesky}}~
\href{mailto:?subject=Editorial&body=https%3A%2F%2Fwww.dropbox.com%2Fscl%2Ffi%2From8ermbyet1p606ms6y0%2Fl-01-EDITORIAL.pdf%3Frlkey%3Dx341kromqoa6a1fyx499z5yj2%26st%3Dmknhry7g%26dl%3D1}{\includesvg[height=0.9em]{/home/alberto/.gbpublisher/fonts/envelope}}~
\href{https://www.facebook.com/sharer/sharer.php?u=https%3A%2F%2Fwww.dropbox.com%2Fscl%2Ffi%2From8ermbyet1p606ms6y0%2Fl-01-EDITORIAL.pdf%3Frlkey%3Dx341kromqoa6a1fyx499z5yj2%26st%3Dmknhry7g%26dl%3D1}{\includesvg[height=0.9em]{/home/alberto/.gbpublisher/fonts/facebook}}~
\href{https://www.linkedin.com/sharing/share-offsite/?url=https%3A%2F%2Fwww.dropbox.com%2Fscl%2Ffi%2From8ermbyet1p606ms6y0%2Fl-01-EDITORIAL.pdf%3Frlkey%3Dx341kromqoa6a1fyx499z5yj2%26st%3Dmknhry7g%26dl%3D1}{\includesvg[height=0.9em]{/home/alberto/.gbpublisher/fonts/linkedin}}~
\href{https://www.reddit.com/submit?url=https%3A%2F%2Fwww.dropbox.com%2Fscl%2Ffi%2From8ermbyet1p606ms6y0%2Fl-01-EDITORIAL.pdf%3Frlkey%3Dx341kromqoa6a1fyx499z5yj2%26st%3Dmknhry7g%26dl%3D1&title=Editorial}{\includesvg[height=0.9em]{/home/alberto/.gbpublisher/fonts/reddit}}~
\href{https://www.researchgate.net/search?q=Editorial}{\includesvg[height=0.9em]{/home/alberto/.gbpublisher/fonts/researchgate}}~
\href{https://t.me/share/url?url=https%3A%2F%2Fwww.dropbox.com%2Fscl%2Ffi%2From8ermbyet1p606ms6y0%2Fl-01-EDITORIAL.pdf%3Frlkey%3Dx341kromqoa6a1fyx499z5yj2%26st%3Dmknhry7g%26dl%3D1&text=Editorial}{\includesvg[height=0.9em]{/home/alberto/.gbpublisher/fonts/telegram}}~
\href{https://wa.me/?text=https%3A%2F%2Fwww.dropbox.com%2Fscl%2Ffi%2From8ermbyet1p606ms6y0%2Fl-01-EDITORIAL.pdf%3Frlkey%3Dx341kromqoa6a1fyx499z5yj2%26st%3Dmknhry7g%26dl%3D1}{\includesvg[height=0.9em]{/home/alberto/.gbpublisher/fonts/whatsapp}}~
\href{https://twitter.com/intent/tweet?url=https%3A%2F%2Fwww.dropbox.com%2Fscl%2Ffi%2From8ermbyet1p606ms6y0%2Fl-01-EDITORIAL.pdf%3Frlkey%3Dx341kromqoa6a1fyx499z5yj2%26st%3Dmknhry7g%26dl%3D1&text=Editorial}{\includesvg[height=0.9em]{/home/alberto/.gbpublisher/fonts/x-twitter}}
}\par
\end{minipage}%
}% FIN GRUPO leftskip
\end{textblock*}


\bigskip
\noindent\rule{\linewidth}{0.4pt}
\bigskip

Para el SENA, Centro de Diseño e Innovación Tecnológica Industrial es un placer presentar la 9na edición de la revista TEINNOVA – Una Mirada de Innovación. Desde la maravillosa región del Eje Cafetero, en Colombia, se acentúa un espacio académico para el encuentro y la reflexión, alrededor de los procesos investigativos.

En esta edición el lector encontrará diversas temáticas, desde bioeconomía y territorio, modelos de innovación para la comunicación organizacional, hasta aspectos de logística inversa propuestos por diferentes autores nacionales e internacionales.

Esta edición buscó reunir estrategias y nuevas tendencias mundiales, en el uso de las tecnologías disruptivas, bioeconomía y territorio con un enfoque hacia la sostenibilidad, presentar el desarrollo de tecnologías y procesos innovadores para la transformación de recursos, así como en la implementación de nuevas técnicas de utilización de maderas plásticas, diseño de prototipos biomédicos y gestión eficiente del uso de recursos naturales y la implementación de tecnologías limpias.

Agradecemos a todos los autores por haber puesto en consideración sus trabajos y plasmar los avances, desarrollo y evolución, de igual forma un especial agradecimiento a todos los pares ciegos que hacen posible que se garantice la rigurosidad escritural en el proceso de gestión editorial. Sin nada más por añadir, esperamos que esta lectura les motive a seguir aprendiendo.

¡A disfrutar cada palabra!

